\chapter{Of Housing and the Floors}
\label{art:housing}

\articlenote{In which is set down the keeping of living quarters for every citizen, the assignment of floors by the silo's zones, the accountability of citizens for their dwelling places, and the principle that every citizen shall have shelter and no citizen shall have more than is needful.}

\section{The Charge of Housing and the Distribution of Quarters}

Every citizen is entitled to a dwelling place appropriate to the citizen's household size and the silo's available stores, held under the administration of the Department of Supply under Article 11. No citizen may claim a dwelling as permanent property, nor may any citizen occupy a floor or dwelling without the approval and assignment of the Head of Supply, and no assignment may be moved or changed without the same approval given in good faith.

\begin{clauses}
\item A dwelling shall be assigned to a citizen household by size and need: a single citizen, an aged citizen living alone, or a child too young for independent keeping shall receive a single room or small quarters; a partnership of two citizens contracted under Article 3 shall receive quarters adequate to the partnership's needs; a parent with one or more children shall receive quarters adequate to the household, with additional space for each child born under the lottery of Article 3.
\item A citizen assigned to a dwelling may make reasonable modifications to the interior (paint, arrangement of common furnishings, minor repairs), provided the citizen restores the dwelling to its original condition upon leaving or ceding the assignment; major structural repairs shall be made only by the Maintenance workers of Supply under Article 11.
\item The floor upon which a dwelling is located shall be occupied primarily by citizens of the trades assigned to that floor under Article 13, and by the households of such citizens, that the silo's work not be endangered by dispersal of the trades across all floors.
\item A citizen's assignment to a dwelling may be revoked by the Head of Supply for failure to keep the dwelling in good repair, for violation of the Pact's requirements concerning the use of a dwelling under Section 3, or for the citizen's reassignment to a new trade requiring residence on a different floor; revocation shall be given in writing, and a citizen may appeal such revocation to Judicial for a ruling.
\end{clauses}

\section{The Zones and the Division of the Silo}

The silo is divided into three zones by the character of the work done there and the citizens assigned to each, exactly as Article 1 sets them down: the Up Top, its uppermost levels, where the Mayor, Judicial, the Sheriff, the teachings of the young, and the common areas of assembly and congregation are kept; the Mid Thirties and the Mids beyond them, where most citizen households are located, and the markets, the Infirmary, Information Technology, and the ordinary trades of the silo are carried out; and the Down Deep, the silo's lowest levels, where the Department of Mechanical, the Department of Mines, the generator, the water treatment and stores, and the deep-storage of the silo's reserves are kept, and where only the citizens assigned to those Departments shall ordinarily dwell.

\begin{clauses}
\item A citizen may not occupy a dwelling in the Down Deep save where that citizen is assigned by trade to the Department of Mechanical, the Department of Mines, or to a support trade kept deep for the works of those Departments, and no other citizen, save by the explicit written permission of the Head of Supply, may lodge or sleep in the Down Deep overnight.
\item The Up Top is reserved for the officers and works of the silo's governance, the teaching of the young, and the common gathering of the Assembly under Article 20; housing in the Up Top is assigned by the Mayor in consultation with the Head of Supply, and is ordinarily provided only to citizens holding offices under Article 4 and the offices of Judicial under Article 6, and to such teachers and scribes as Article 15 and Article 21 require.
\item A citizen of the Mids may petition to work in the Up Top or the Down Deep, and such petition shall be considered under Article 13's labor-assignment process; a citizen successful in such a petition shall be assigned housing on the floor where the citizen's work is situated, or given permission to continue lodging in the Mids and to travel the stairs daily.
\end{clauses}

\section{The Keeping of Dwelling Quarters}

A citizen occupying a dwelling is answerable for the condition of that dwelling and for the conduct of the citizen's household within it. The Pact does not compel citizen-to-citizen courtesy or affection; it does compel that a citizen's household not become a hazard or a breach to other citizens in the dwelling or the floor.

\begin{clauses}
\item A citizen's household shall not produce vermin, filth, or odor that spreads beyond the dwelling itself, nor shall a citizen's household contain more citizens than the dwelling is assigned to accommodate, nor shall a citizen harbor a citizen unknown to the rolls or hide a child born without the lottery's sanction.
\item A citizen may not keep animals not approved under Article 11's Section 5 (sanctioned animals), nor forbid entry to the Maintenance workers of Supply where a repair to the dwelling's systems is required.
\item A citizen may not modify the structural integrity of a dwelling, nor seal off doors or passages, nor make permanent alterations without the approval of the Head of Supply and the approval, where relevant, of the Head of Mechanical for concerns affecting the silo's structure.
\item A citizen who breaches this section's requirements is answerable under Article 17; if the breach endangers the floor's stability or the water, air, or waste-flow systems under Mechanical's charge, it is among the gravest offenses under Article 17.
\end{clauses}

\section{The Privacy of Dwelling Quarters and the Limits of Search}

The Pact recognizes that a citizen's dwelling is a citizen's refuge, and that some matters of household life are private. Nevertheless, the silo's safety and the enforcement of the Pact take precedence over dwelling privacy where the Pact judges they must.

\begin{clauses}
\item The Sheriff may enter a dwelling to stop a crime in progress, to investigate a death, or to arrest a citizen for an offense. In all other cases, the Sheriff shall request entry, and may enter only if a citizen of the household consents or if the Judge has ordered a search under Article 6.
\item The Officers of Judicial, when executing a search ordered by the Judge under Article 6, may search a dwelling and the citizen's private possessions; such a search shall be conducted in daylight or by lamp, shall not destroy what is searched unnecessarily, and shall be reported to Judicial for the archive under Article 6.
\item A search conducted without a Judge's order, or conducted by an officer without the Judge's written order, is a violation of the Pact, and a citizen may refuse entry and may file a complaint with Judicial.
\end{clauses}

\section{Housing and the Equality of Citizens}

The Pact provides that every citizen shall have shelter adequate to the citizen's household, and that no citizen shall go unsheltered or displaced without cause grounded in the Pact itself. The Head of Supply, in consultation with the Mayor, shall ensure that no citizen who fulfills the duties of the Pact is left without a dwelling.

\begin{clauses}
\item Where a citizen is displaced from a dwelling by catastrophe, by the citizen's release from the mines or from cleaning, or by reassignment of the citizen's trade to a different floor, the Head of Supply shall assign new quarters without undue delay.
\item A citizen facing removal from a dwelling for breach of this Article's requirements may petition Judicial for a ruling; Judicial may find the breach overstated, may impose a deadline for remedy before removal takes effect, or may affirm the removal and order the citizen to a temporary holding pending reassignment.
\item The silo's citizens are entitled to housing, not to choice of housing; a citizen may petition for a change of quarters, and the Head of Supply shall consider such a petition in good faith, but is not bound by a citizen's preference where silo needs require otherwise.
\end{clauses}
